\scanpage{302}
\mspage{26}
% archive sheet number: 301

i.e. \unsure{toute|tout} \unsure{flèche|flèches} $f$ de $\mathrm{Cat}$ se factorise en $f=pi$, où $i\in\mathrm{Cof}_{\unsure{W|w}}$ et $p\in W$. (Cela signifie aussi que $\mathrm{Cat}$ munie de $(W,\mathrm{Cof}_W)$ est une catégorie à cofibrations de K. S. Brown--Anderson, et bien des résultats de la théorie des \unsure{Brown-Quillen|Brown-oreilles} sont applicables\deleted{\unsure{.)}}. En particulier, on a une bonne théorie des carrés $W$-cocartésiens.)

\uline{$\mathrm{Loc}\,6\,\mathrm B$} Si \deleted{$\unsure{X}$} dans $\mathrm{Cat}$ est réunion de deux ouverts $U,V$, d'intersection $w$, alors le carré
% slanted
\margin{En fait, $\mathrm{Loc}\,6\,\mathrm B$ est conséquence de $\mathrm{Loc}\,6\,\mathrm A$, \deleted{\unsure{(qui} je prouve} \inserted{laquelle lui} est \uline{équivalente} à $\mathrm{Loc}\,6a$ ? Cf esquisse de démonstration \deleted{page} des 4.4 bis pages 66--68.}
\diagram{\begin{tikzcd}[ampersand replacement=\&]
w\arrow[r]\arrow[d]\&U\arrow[d]\\
V\arrow[r,dash]\&X
\end{tikzcd}}
est $W$-cocartésien.

D'autre part, la condition $\mathrm{Loc}\,6\,\mathrm A$ \deleted{\unsure{est cette}} \inserted{peut être} \deleted{équivalente à} \inserted{remplacée par} la condition (en apparence plus faible)

\uline{$\mathrm{Loc}\,6\,\mathrm A'$} Tout morphisme de Dwyer est dans $\mathrm{Cof}_W$.

\uline{Corollaire} 3.7. La condition $\mathrm{Loc}\,6\deleted{\unsure{(ab)}}$ \inserted{i.e $\mathrm{Loc}\,6(a,b)$} (ou encore la condition équivalente $\mathrm{Loc}\,6'$) équivaut à la conjonction des conditions $\mathrm{Loc}\,6\,\mathrm A$ \inserted{(ou $\mathrm{Loc}\,6\,\mathrm A'$)} et $\mathrm{Loc}\,6\,\mathrm B$ ci-dessus, et de la condition

\uline{$\mathrm{Loc}\,6\,\mathrm C$} : Pour \deleted{tout} carré \inserted{$W$-}cocartésien
\diagram{\begin{tikzcd}[ampersand replacement=\&]
X_0\arrow[r,"\unsure{i|j}"]\arrow[d]\&X_1\arrow[d]\\
X_{\unsure{2|1}}\arrow[r,"j"]\&\unsure{X|X^{\prime}}
\end{tikzcd}}
dans $\mathrm{Cat}$, $j\in W\Longrightarrow i\in W$.
